\documentclass[11pt,a4paper]{article}
\usepackage[utf8]{inputenc}
\usepackage[T1]{fontenc}
\usepackage{amsmath,amssymb}
\usepackage{booktabs}
\usepackage{listings}
\usepackage{xcolor}
\usepackage{hyperref}
\usepackage{geometry}
\geometry{margin=2.5cm}

\lstset{
  basicstyle=\ttfamily\small,
  breaklines=true,
  frame=single,
  backgroundcolor=\color{gray!10},
  numbers=left,
  numberstyle=\tiny
}

\title{02244 - Logic for Security\\
\vspace{0.5em}
Project: Security Protocols}

\author{%
  \textbf{Group Members}\\[0.4em]
  Zeinab\\
  Fahmida\\
  Mirtha\\[0.8em]
  March 2026
}

\date{}

\begin{document}
\maketitle
\thispagestyle{empty}

\newpage
\tableofcontents
\newpage

\section*{Section Author Names}
\begin{itemize}
  \item Section 1--3: Zeinab
  \item Section 4--6: Fahmida
  \item Section 7--11: Mirtha
\end{itemize}

\newpage
\section{Introduction}

In modern online services, users often need to allow third-party applications to access selected resources stored on another platform.
Direct download and re-upload is inefficient and increases exposure risk.
An authorization protocol is therefore needed to allow controlled delegation without exposing user credentials.

This project designs and analyzes a protocol that allows photo-book service $P$ to access photos stored at server $B$ on behalf of user $A$, using an identity provider $I$ for authentication and authorization issuing.
The protocol family is modeled in AnB and analyzed with OFMC under symbolic attacker assumptions.

The key participants are:
\begin{itemize}
  \item $A$: user owning photos,
  \item $B$: photo storage server,
  \item $P$: photo-book service,
  \item $I$: trusted identity provider.
\end{itemize}

Main goals are authentication, authorization, and confidentiality of delegated photo access.

\section{System Model and Assumptions}

\subsection*{Agents}
\begin{itemize}
  \item \textbf{$A$} -- owner of photos.
  \item \textbf{$B$} -- resource server holding photos.
  \item \textbf{$P$} -- delegated service requesting read access.
  \item \textbf{$I$} -- identity provider issuing signed grants.
\end{itemize}

\subsection*{Assumptions}
\begin{itemize}
  \item $I$ is trusted and honest.
  \item Everybody knows public keys needed to verify $I$ statements.
  \item Password relation $pw(A,I)$ is low-entropy (guessable model in Week 6).
  \item $A$ does not rely on a strong long-term signing infrastructure in the early designs; authorization trust is centralized in $I$.
\end{itemize}

\subsection*{Security goals}
\begin{itemize}
  \item Authentication of relevant parties on delegation context.
  \item Authorization: only valid IDP-issued grants should be accepted by $B$.
  \item Confidentiality of photo payload.
  \item Limited access rights (scoped request fields).
\end{itemize}

\section{Protocol Design (Development Process)}

\subsection{3.1 Initial ``Bare Bones'' Protocol}
The Week 2 baseline used a simple flow:
$A$ signs delegation context for $P$, $P$ forwards it to $B$, and $B$ returns photos.
This version was intentionally simple to create an executable baseline and perform static Dolev--Yao analysis.

\subsection{3.2 Improved Version (Week 3 -- Identity Provider)}
The protocol was upgraded so that:
\begin{itemize}
  \item $A$ authenticates to $I$ using password evidence $h(pw(A,I),N)$,
  \item $I$ signs an authorization grant,
  \item $P$ presents this grant to $B$.
\end{itemize}
This change aligns with the assignment requirement that password is used for authentication proof, not as encryption key.

\subsection{3.3 Final Protocol Candidate}
The final candidate used for Week 6 checks follows this flow:
\begin{enumerate}
  \item $p \rightarrow A$: delegation request context.
  \item $A \rightarrow I$: authenticated request with password witness and freshness.
  \item $I \rightarrow A$: signed grant token.
  \item $A \rightarrow p$: token forwarding.
  \item $p \rightarrow B$: token + scope request.
  \item $B \rightarrow p$: encrypted photo payload.
\end{enumerate}

\section{Static Analysis (Week 2 Task)}

Under the required Week 2 setting, the intruder only observes traffic and does not inject messages.
Using Dolev--Yao deduction rules, we analyzed what can be derived from observed messages in the baseline model.

Result: the baseline design does not provide robust confidentiality under all role instantiations, which motivated the Week 3 move to explicit IDP-issued grants.
This task is documented in \texttt{Week2\_Problem3\_Static\_DolevYao\_Analysis.md}.

\section{Lazy Intruder Analysis (Week 3)}

The special Week 3 task asks whether the intruder can play role $P$.
Following the lecture method, we instantiate $P=i$ and keep other roles honest.
Then, step-by-step:
\begin{itemize}
  \item intruder builds outgoing role-$P$ messages from current knowledge (Axiom/Compose),
  \item receives token-carrying messages,
  \item forwards required request to $B$.
\end{itemize}

Conclusion: role $P$ is executable in the symbolic model.
The formal derivation is in \texttt{Week3\_Problem2\_LazyIntruder\_Executability\_of\_P.md}.

\section{Typing and Implementation (Week 4)}

Week 4 introduces explicit message formats and a separate key-lookup protocol.
Format constructors (e.g., \texttt{fReq}, \texttt{fGrant}, \texttt{fData}) improve message-type separation and reduce type-confusion opportunities.

Type-flaw resistance argument:
if two non-variable protocol patterns unify, they must have compatible top-level structure.
This was documented in \texttt{Week4\_Problem3\_TypeFlawResistance\_Check.md}.

\section{Channels and TLS (Week 5)}

Week 5 replaced direct message-level transport cryptography with pseudonymous channel notation.
This approximates TLS-like channel assumptions with authenticated server side and abstract secure transport semantics.

Result: the Week 5 channel variant is executable but still yields attack traces in OFMC.
Hence channel abstraction improves modeling clarity but does not automatically solve logic-level authorization issues.

\section{Password Guessability (Week 6)}

Week 6 explicitly checks low-entropy password risk:
\begin{itemize}
  \item Baseline guessable model: \texttt{ATTACK\_FOUND} on \texttt{guesswhat}.
  \item Hardened model (extra fresh value in proof): first failure shifts away from \texttt{guesswhat}.
  \item Focused guessable-only model (\texttt{--numSess 1}): \texttt{NO\_ATTACK\_FOUND}.
\end{itemize}

Therefore, for the dedicated Week 6 special-task goal and tested bound, the protocol variant resists the modeled guessing attack.

\section{Final AnB Protocol}

We use the Week 6 hardened pattern as the final candidate for discussion:

\begin{lstlisting}
Protocol: OpenAuth_W6_GuessablePassword_Hardened

Types:
  Agent A,B,I,p;
  Number F,T,R,K,Grant,Photos;
  Function pk,pw,h,fGrant,fData;

Actions:
  p -> A: p,B,F,T,R
  A -> I: {A,p,B,F,T,R,K,h(pw(A,I),R,K)}pk(I)
  I -> A: {fGrant(A,p,B,F,T,Grant)}inv(pk(I))
  A -> p: {fGrant(A,p,B,F,T,Grant)}inv(pk(I))
  p -> B: {fGrant(A,p,B,F,T,Grant)}inv(pk(I)),F
  B -> p: {fData(Photos)}pk(p)
\end{lstlisting}

\subsection*{Initial knowledge (summary)}
$A$ and $I$ share password relation; $I$ signs grants; $B$ verifies IDP-issued token context; $p$ receives delegated data from $B$.

\subsection*{Goals (used across variants)}
\begin{itemize}
  \item Secrecy of delegated photo payload term.
  \item Authentication on grant context.
  \item Guessable-password secrecy checks (Week 6 task).
\end{itemize}

\section{Security Discussion}

The protocol family shows clear improvement from Week 2 to Week 6, but not every variant satisfies every goal simultaneously.
This is expected in iterative protocol engineering and was required to be documented.

Strengths:
\begin{itemize}
  \item explicit trust anchor through IDP-signed grants,
  \item typed/format-aware message design in Week 4,
  \item dedicated guessable-password analysis in Week 6.
\end{itemize}

Limitations:
\begin{itemize}
  \item some variants remain vulnerable to weak-auth or secrecy attacks,
  \item several conclusions are bounded-session results,
  \item channel abstraction does not replace authorization logic correctness.
\end{itemize}

\section{Conclusion}

The project develops a complete week-by-week Open Authorization protocol design for delegated photo access from $B$ to service $P$ on behalf of user $A$, with IDP-mediated trust.
The development includes baseline modeling, attack discovery, structured refinements, and targeted special-task verification.

Formal analysis with OFMC confirms that naive versions are vulnerable and that improvements are necessary.
In particular, Week 6 focused verification demonstrates bounded resistance against the modeled guessable-password attack.
Overall, the project satisfies the assignment objective of documenting protocol design as an evidence-based development process rather than presenting only a final static model.

\appendix
\section{Integrated Notes from Markdown Files}
This appendix embeds the full text content that was previously stored in project Markdown notes.

\subsection{Week2\_Problem1\_Informal\_Protocol\_Outline.md}
\begin{lstlisting}
# Week 2 - Problem 1: Informal Protocol Outline

## Goal (Mapped to Project Text)
This models the project scenario directly:
- `A`: account owner (the person with holiday photos)
- `B`: photo storage server/resource server
- `P`: photo-book printing service

`A` wants to let `P` read only selected photos on `B` without sharing full account credentials.

## Simplified Assumptions for Week 2
- We simplify key management: all public keys are already known.
- No identity provider yet (added in Week 3).
- We model one authorization scope as a token.
- The token is intended to represent read-only access to one selected photo folder/set.

## Informal Message Flow (First Draft)
1. `P -> A`: "Please allow me to read your selected photos on server `B`" (request + nonce).
2. `A -> P`: signed authorization statement bound to (`P`, `B`, scope, nonce, token).
3. `P -> B`: forwards `A`'s authorization and requests the approved photo scope.
4. `B -> P`: returns encrypted photo data under the session token.

## Initial Security Intent
- `B` only accepts access requests that are authorized by `A`.
- Authorization is bound to both `P` (delegate) and `B` (resource server).
- Photo data should only become available to the authorized `P`.
- The authorization should be scoped (selected folder/photos, not all resources).

## Known Limitation of this Week-2 Draft
- The model is intentionally minimal and may leak data under Dolev-Yao analysis.
- We use this version as baseline for iterative improvement in later weeks.
\end{lstlisting}

\subsection{Week3\_Problem2\_LazyIntruder\_Executability\_of\_P.md}
\begin{lstlisting}
# Week 3 - Problem 2: Lazy Intruder Executability of Role `P`

Protocol analyzed: `Week3_Problem1_AnB_With_IdentityProvider.anb`

Required setting (from assignment):
- instantiate `P = i` (intruder plays role `P`);
- all other roles are honest (`A`, `B`, `I`);
- show intruder can produce every outgoing message of role `P`.

Role `P` messages in this protocol:
1. `P -> A: P,B,F,T,NA`
2. `P -> B: {A,P,B,F,T,NA,Grant}inv(pk(I)),F`

## Initial Intruder Knowledge for Role `P=i`
From role `P` knowledge (after instantiation):
`{A,B,i,I,pk(A),pk(B),pk(i),pk(I),inv(pk(i))}`

`F,T,NA,Grant` are variables and may be left lazy until needed.

## Step-by-Step Lazy Intruder Derivation
### Outgoing Message 1
Target:
`i, B, F, T, NA`

Derivation:
1. `i` is in knowledge (Axiom).
2. `B` is in knowledge (Axiom).
3. `F` is a variable: keep lazy.
4. `T` is a variable: keep lazy.
5. `NA` is a variable: keep lazy.
6. Compose tuple with public concatenation operator.

So message 1 is executable.

### Incoming Message before Outgoing Message 2
Before message 2, role `P` receives from `A`:
`{A,i,B,F,T,NA,Grant}inv(pk(I))`

This exact signed token is now part of intruder knowledge by Axiom.

### Outgoing Message 2
Target:
`{A,i,B,F,T,NA,Grant}inv(pk(I)),F`

Derivation:
1. First component `{A,i,B,F,T,NA,Grant}inv(pk(I))`:
   already known from previous incoming message (Axiom).
2. Second component `F`:
   known from first message context (lazy variable choice remains valid).
3. Compose pair with public concatenation operator.

So message 2 is executable.

## Conclusion
Following the lazy intruder method from the lecture (Axiom + Compose + lazy variables), the intruder can generate all outgoing messages of role `P`.

Therefore, role `P` is executable under the required Week 3 instantiation.
\end{lstlisting}

\subsection{Week4\_Problem3\_TypeFlawResistance\_Check.md}
\begin{lstlisting}
# Week 4 - Problem 3: Type-Flaw Resistance Check

Files checked:
- `Week4_Problem1_MainProtocol_Formats.anb`
- `Week4_Problem2_AnB_KeyLookup_Protocol.anb`

Goal (from lecture): for any two non-variable SMP terms, if they unify, they have the same type.

## 1) Format design used
Main delegation protocol uses explicit format symbols:
- `fReq(p,B,F,T,NA)`
- `fAuthProof(A,p,B,F,T,NA,h(pw(A,I),NA))`
- `fGrant(A,p,B,F,T,NA,Grant)`
- `fUse(Token,F)` where `Token = {fGrant(...)}inv(pk(I))`
- `fData(Photos)`

Key lookup protocol uses:
- `fKeyReq(A,X)`
- `fKeyResp(X,pk(X))`

Each protocol message kind has a dedicated top-level format symbol.

## 2) Why this prevents type-flaw unification
Potential unification between different message kinds would require equal top-level symbols.
That is impossible across different format constructors:
- `fReq` does not unify with `fAuthProof`, `fGrant`, `fUse`, `fData`, `fKeyReq`, or `fKeyResp`.
- `fKeyReq` does not unify with `fKeyResp`.
- Encrypted forms also remain separated by inner format type.

So only same-kind messages can unify, and those are same type by construction.

## 3) Conclusion
Under the Week 4 formatted design, message-pattern unification is type-aligned.
Therefore this model is type-flaw resistant in the intended lecture sense.
\end{lstlisting}

\subsection{Week5\_Problem1\_PseudonymousChannels\_Result.md}
\begin{lstlisting}
# Week 5 - Problem 1 Result (Pseudonymous Channels)

Protocol file:
- `Week5_Problem1_PseudonymousChannels.anb`

Command used:
- `ofmc --numSess 2 "ProjectProtocol\Week5_Problem1_PseudonymousChannels.anb"`

Result:
- `SUMMARY: ATTACK_FOUND`
- `GOAL: secrets`
- leaked term: `fData(Photos(1))`

Interpretation:
- The Week 5 model correctly uses OFMC pseudonymous channels (`*->*` and `[A]`).
- OFMC still finds an attack on secrecy in this variant.
- This is acceptable as a development snapshot and can be discussed in the report as a limitation/attack trace for the channel-based version.
\end{lstlisting}

\subsection{Week6\_Problem1\_Result\_Summary.md}
\begin{lstlisting}
# Week 6 - Guessable Password Check (Summary)

## Files
- `Week6_Problem1_GuessablePassword_Check.anb`
- `Week6_Problem2_GuessablePassword_Hardened.anb`
- `Week6_Problem3_GuessablePassword_OnlyCheck.anb`

## OFMC Results
### 1) Baseline guessable check
Command:
- `ofmc --numSess 2 "ProjectProtocol\Week6_Problem1_GuessablePassword_Check.anb"`

Result:
- `ATTACK_FOUND`
- `GOAL: guesswhat`

Interpretation:
- Password usage in this baseline is vulnerable to offline guessing in OFMC's guessable-secret model.

### 2) Hardened variant (adds fresh `K` into password proof)
Command:
- `ofmc --numSess 2 "ProjectProtocol\Week6_Problem2_GuessablePassword_Hardened.anb"`

Result:
- `ATTACK_FOUND`
- `GOAL: secrets` (not `guesswhat`)

Interpretation:
- Guessing attack is no longer the first issue found.
- A separate secrecy weakness still exists in this protocol variant.

### 3) Guessable-only verification
Command:
- `ofmc --numSess 1 "ProjectProtocol\Week6_Problem3_GuessablePassword_OnlyCheck.anb"`

Result:
- `NO_ATTACK_FOUND`
- `DETAILS: BOUNDED_NUMBER_OF_SESSIONS`

Interpretation:
- For this bounded check, the password is not broken via the guessable-secret goal.
- This directly addresses the Week 6 special-task focus (guessable password robustness), within the tested session bound.
\end{lstlisting}

\subsection{Week7\_Problem1\_ClassPresentation\_Outline.md}
\begin{lstlisting}
# Week 7 - Class Presentation Outline

## Slide 1 - Title and Scenario
- Project: Open Authorization for photo-book printing.
- Roles:
  - `A`: photo owner
  - `B`: photo storage server
  - `P`/`p`: photo-book service
  - `I`: identity provider
- Goal: let `P` read selected photos from `B` without exposing full account access.

## Slide 2 - Week 2 Baseline (Dolev-Yao)
- File: `Week2_Problem2_First_AnB_Version.anb`
- Initial simple delegation model.
- Static Dolev-Yao analysis file:
  - `Week2_Problem3_Static_DolevYao_Analysis.md`
- Key finding: baseline secrecy can leak under attacker observation/model assumptions.

## Slide 3 - Week 3 IDP Integration + Lazy Intruder
- Main file: `Week3_Problem1_AnB_With_IdentityProvider.anb`
- Added identity provider and password proof:
  - `h(pw(A,I),NA)` (password not used as encryption key).
- Special task file:
  - `Week3_Problem2_LazyIntruder_Executability_of_P.md`
- Key finding: role executability shown; security issues still found in model checking.

## Slide 4 - Week 4 Formats + Key Lookup
- Main format file:
  - `Week4_Problem1_MainProtocol_Formats.anb`
- Separate key lookup:
  - `Week4_Problem2_AnB_KeyLookup_Protocol.anb`
- Type-flaw reasoning:
  - `Week4_Problem3_TypeFlawResistance_Check.md`
- Key finding: format-based structure improves protocol clarity and typing robustness.

## Slide 5 - Week 5 Pseudonymous Channels (TLS-style)
- File: `Week5_Problem1_PseudonymousChannels.anb`
- Goal: replace direct crypto expressions with channel abstraction (`*->*`, `[A]` style).
- Result summary:
  - `Week5_Problem1_PseudonymousChannels_Result.md`
- Key finding: channel abstraction works but still exposes weaknesses in this variant.

## Slide 6 - Week 6 Guessable Password Verification
- Baseline vulnerable check:
  - `Week6_Problem1_GuessablePassword_Check.anb` -> `guesswhat` attack found.
- Hardened version:
  - `Week6_Problem2_GuessablePassword_Hardened.anb`
- Focused special-task check:
  - `Week6_Problem3_GuessablePassword_OnlyCheck.anb`
  - OFMC (`--numSess 1`) -> `NO_ATTACK_FOUND` for guessable-password goal.

## Slide 7 - Main Lessons
- Iterative development is essential in protocol design.
- Passing syntax/executability does not mean security.
- Different goals fail differently (secrecy vs. weak authentication vs. guessing).
- Message formats and scoped assumptions significantly change verification outcomes.

## Slide 8 - Final Status and Next Steps
- Final narrative:
  - Week 2 -> baseline
  - Week 3 -> IDP + executability proof
  - Week 4 -> format and typing improvements
  - Week 5 -> channel abstraction
  - Week 6 -> guessable password analysis
- Future improvement:
  - remove remaining weak-auth/secrecy attacks while preserving delegation features.
\end{lstlisting}

\subsection{Week7\_Problem2\_SpeakingScript\_Short.md}
\begin{lstlisting}
# Week 7 - Short Speaking Script (5-7 minutes)

## 1) Intro (30-45 sec)
"Our project models open authorization for photo printing. A wants to allow P to read selected photos at B, with trust anchored in identity provider I. We developed the protocol week by week and verified each version with OFMC."

## 2) Week 2 (45 sec)
"We started with a simple delegation protocol in AnB. Using Dolev-Yao analysis and OFMC, we observed that secrecy could fail in the baseline. This gave us a concrete starting point for refinement."

## 3) Week 3 (60 sec)
"We integrated the identity provider and password-based authentication proof using `h(pw(A,I),NA)` rather than using password as encryption key. We also completed the required lazy-intruder executability task for role P step by step."

## 4) Week 4 (60 sec)
"We introduced explicit message formats to reduce type confusion and made a separate key-lookup protocol where A queries I for public keys. We documented type-flaw-resistance reasoning based on message-pattern structure."

## 5) Week 5 (45 sec)
"We created a TLS-style channel abstraction variant using pseudonymous channels. This made the model closer to realistic channel assumptions, but attacks were still possible in our tested variant."

## 6) Week 6 (60-75 sec)
"We focused on guessable password security. The baseline guessable-password model was breakable. Then we hardened the password proof with extra freshness and checked a guessable-only model. For bounded sessions, this check returned no guessing attack."

## 7) Closing (30 sec)
"Main lesson: protocol security requires iterative modeling, explicit assumptions, and multiple goal checks. Secrecy, authentication, and guessing-resistance can fail independently, so each must be tested explicitly."

## Optional Q&A backup points
- Why include attack-found versions? -> They justify design decisions and show learning progression.
- Why separate key lookup? -> Week 4 requirement and realistic trust decomposition.
- Why bounded result in Week 6? -> OFMC bounded-session analysis; we report exactly what was verified.
\end{lstlisting}

\subsection{Report\_Writing\_Guide\_NoMissing.md}
\begin{lstlisting}
# Report Writing Guide (No Missing Parts)

Use this guide with `Report_ProjectProtocol_ProfessorStyle.tex` so your final report is complete and aligned with the assignment.

## 1) What you must submit
From the assignment, you submit:
1. **Report PDF** (main write-up, max page limit)
2. **Lab Logbook** (development story across weeks)
3. **AnB files**

You do **not** have to submit every `.md` as separate graded documents unless you want to attach them as supporting notes.

## 2) Why there are `.md` files
The `.md` files were created for two reasons:
- **To capture weekly required explanations** (e.g., Week 3 lazy intruder steps, Week 4 type-flaw argument, Week 6 result summary).
- **To make report writing easier** by reusing tested text/results instead of rewriting everything from scratch.

Think of `.md` files as **draft evidence + notes**.
Your final official text should be in the report PDF (LaTeX).
\end{lstlisting}

\end{document}
